\begin{problem}[第36届全国中学生物理竞赛决赛][涡轮增压器热力学过程]{114}
    汽车以发动机气缸中的空气与汽油发生燃烧产生动力，发动机的输出功率与进入气缸中的空气质量成正比。为了提高动力，很多汽车加装了涡轮增压器，在空气进入发动机气缸之前对它进行压缩，以增加空气密度；为了进一步提高空气密度，还通过一个中间冷却器，使空气温度降低后再进入气缸。在一个典型的装置中，进入涡轮增压器的初始空气压强为一个大气压 $p_{0} = 1.01 \times 10^{5} \,\mathrm{Pa}$，温度 $T_{0} = 15^{\circ}\mathrm{C}$；通过涡轮增压器后，空气绝热压缩至 $p_{1} = 1.45 \times 10^{5} \,\mathrm{Pa}$；然后又通过中间冷却器，让空气等压地降至初始温度 $T_{0}$ 并充满气缸。假设空气为理想气体，已知气缸容积 $V_{0} = 400 \,\mathrm{cm}^{3}$，空气摩尔质量 $M = 29.0 \,\mathrm{g} \cdot \mathrm{mol}^{-1}$，摩尔热容比 $\gamma = 7/5$；普适气体常量 $R = 8.31 \,\mathrm{J} \cdot \mathrm{K}^{-1} \cdot \mathrm{mol}^{-1}$。
    \begin{subproblem}
        求此时气缸内空气的质量；与未经增压和冷却的情形相比，发动机的输出功率增加到多少倍？
    \end{subproblem}
    \begin{subproblem}
        如果没有中间冷却器，仅经过涡轮增压器把空气压缩后输入到气缸中，相应的空气质量为多少？与未经增压的情形相比，发动机输出功率增加到多少倍？
    \end{subproblem}
    \begin{subproblem}
        在经过涡轮增压器和中间冷却器的热力学过程中，求外界对进入气缸的气体所做的功、及气缸中的气体在上述两个过程中各自的熵变。
    \end{subproblem}
\end{problem}